% Part: normal-modal-logic
% Chapter: frame-correspondence
% Section: equivalence-S5

\documentclass[../../../include/open-logic-section]{subfiles}

\begin{document}

\olfileid{nml}{frd}{es5}

\olsection{د هم ارزښتۍ اړيکې او \Log{S5}}

مودالي منطق \Log{S5} د هغو !!{formula}s د مجموعې په توګه پېژندل
کېږي چې په ټولو کلي چوکاټونو کښې معتبر وي؛ يعنې هرې نړۍ ته له هرې
نړۍ، د هغې له خپل ځان څخه هم، لاسرسے وي۔ په داسې حالت کښې
$\Box$ له لزوم او $\Diamond$ له امکان سره سمون لري: $\Box !A$
هله رښتيا دے چې $!A$ په \emph{هرې} نړۍ کښې رښتيا وي، او
$\Diamond !A$ هله رښتيا دے چې $!A$ په \emph{کومې} نړۍ کښې
رښتيا وي۔ دا هم څرګندېږي چې \Log{S5} د هغو !!{formula}s په
توګه هم پېژندل کېدلے شي چې په ټولو انعکاسي، متناظرو او انتقالي
چوکاټونو کښې معتبر وي؛ يعنې په هغو چوکاټونو کښې چې د لاسرسي
اړيکه ئې د \emph{هم ارزښتۍ اړيکه} وي۔

\begin{defn}
  پر $W$ دوه‌ځاييزه اړيکه $R$ هله او يوازې هله د
  \emph{هم ارزښتۍ اړيکه} ده چې انعکاسي، متناظره او انتقالي وي۔
  پر $W$ اړيکه $R$ هله او يوازې هله \emph{کلي} ده چې د ټولو
  $u,v \in W$ دپاره $Ruv$ وي۔
\end{defn}

ځکه چې \Ax{T}، \Ax{B} او \Ax{4} په ترتيب سره انعکاسي،
متناظره او انتقالي چوکاټونه پېژني، نو هغه چوکاټونه چې د لاسرسي
اړيکه ئې د هم ارزښتۍ اړيکه وي، کټ مټ هغه دي چې دا درې واړه
!!{formula}s پکې معتبر وي۔ د هم ارزښتۍ اړيکې د !!{formula}s
په نورو ترکيبونو هم پېژندل کېدلے شي، ځکه هغه شرطونه چې پرې مو
د هم ارزښتۍ اړيکې تعريف کړې دي، پر~$R$ د نورو پېژندل شوو
شرطونو له ترکيبونو سره هم‌ارز دي۔

\begin{prop}\ollabel{prop:equivalences}
  لاندې شرطونه له يو بل سره هم‌ارز دي:
  \begin{enumerate}
  \item $R$ د هم ارزښتۍ اړيکه ده؛
  \item $R$ انعکاسي او اقليدسي ده؛
  \item $R$ مسلسل، متناظره او اقليدسي ده؛
  \item $R$ مسلسل، متناظره او انتقالي ده۔
  \end{enumerate}
\end{prop}

\begin{proof}
  تمرين۔
\end{proof}

\begin{prob}
  د لاندې خبرو په ښودلو \olref[nml][frd][es5]{prop:equivalences}
  ثابته کړئ:
  \begin{enumerate}
  \item که $R$ متناظره او انتقالي وي، نو اقليدسي ده۔
  \item که $R$ انعکاسي وي، نو مسلسل ده۔
  \item که $R$ انعکاسي او اقليدسي وي، نو متناظره ده۔
  \item که $R$ متناظره او اقليدسي وي، نو انتقالي ده۔
  \item که $R$ مسلسل، متناظره او انتقالي وي، نو انعکاسي ده۔
  \end{enumerate}
  څرګنده کړئ چې دا خبرې ولې د دغو شرطونو د هم‌ارزۍ د ثبوت
  دپاره بسنه کوي۔
\end{prob}

\olref{prop:equivalences} د \olref[prf][prs]{prop:S5} معنايي
سياله ده: د هغو چوکاټونو مودالي منطق په هم‌ارزه بڼه پېژني چې
په هغو کښې~$R$ د هم ارزښتۍ اړيکه وي؛ دې منطق ته په دوديزه توګه
\Log{S5} وايي۔

د کليو او د هم ارزښتۍ اړيکو ترمنځ څۀ اړيکه ده؟ که څۀ هم هره
کلي اړيکه د هم ارزښتۍ اړيکه ده، خو څرګنده ده چې د هم ارزښتۍ
هره اړيکه کلي نۀ ده۔ بيا هم هغه !!{formula}s چې په ټولو
کلي اړيکو کښې معتبر وي، کټ مټ هماغه دي چې په ټولو د هم
ارزښتۍ اړيکو کښې معتبر وي۔

\begin{prop}
  فرض کړئ $R$ د هم ارزښتۍ اړيکه ده، او د هر $w \in W$ دپاره
  د~$w$ \emph{د هم ارزښتۍ ټولګے} د $[w] = \{w'\in W :
  Rww'\}$ په توګه تعريف کړئ۔ بيا:
  \begin{enumerate}
  \item $w \in [w]$؛
  \item $R$ د هم ارزښتۍ پر هر ټولګي $[w]$ کلي ده؛
  \item د هم ارزښتۍ د ټولګيو مجموعه~$W$ په داسې فرعي سټونو
    وېشي چې يو له بل څخه بېل دي او په ګډه ټول سټ رانغاړي۔
  \end{enumerate}
\end{prop}

\begin{prop}\ollabel{prop:S5=univ}
  !!^a{formula}~$!A$ په ټولو هغو چوکاټونو
  $\mModel{F} =\tuple{W,R}$ کښې معتبر دے چې~$R$ پکې د هم
  ارزښتۍ اړيکه وي، هله او يوازې هله چې په ټولو هغو چوکاټونو
  $\mModel{F} =\tuple{W,R}$ کښې معتبر وي چې~$R$ پکې کلي
  وي۔ نو د کلي چوکاټونو منطق هماغه \Log{S5} دے۔
\end{prop}

\begin{proof}
  نېغ په نېغه ښکاري چې پر~$W$ هره کلي اړيکه~$R$ د هم
  ارزښتۍ اړيکه ده۔ نو که~$!A$ په ټولو هغو چوکاټونو کښې
  معتبر وي چې~$R$ پکې د هم ارزښتۍ اړيکه وي، په ټولو کلي
  چوکاټونو کښې هم معتبر دے۔ د بل لوري دپاره د دعوې مقابل
  نقيض ثابتوو: فرض کړئ $!B$ داسې !!a{formula} وي چې په
  يوه نړۍ~$w$ کښې، په داسې مدل
  $\mModel{M} = \tuple{W, R, V}$ کښې نارښتيا وي چې پر
  چوکاټ~$\tuple{W,R}$ ولاړ دے او~$R$ پر~$W$ د هم ارزښتۍ
  اړيکه ده۔ نو $\mSat/{M}{!B}[w]$۔ يو مدل
  $\mModel{M}' = \tuple{W',
    R', V'}$ په دې ډول تعريف کړئ:
  \begin{enumerate}
  \item $W' = [w]$؛
  \item $R'$ پر $W'$ کلي ده؛
  \item $V'(p) = V(p) \cap W'$۔
  \end{enumerate}
  (نو په~$\mModel{M}'$ کښې د نړۍو سټ~$W'$ د
  \olref{fig:partition} په خړ رنګ شوې برخه کښې ښودل شوے دے۔)
  په اسانه ښکاري چې $R$ او $R'$ پر~$W'$ يو شان دي۔ بيا د
  !!{formula}s پر جوړښت په استقرا ښودلے شو چې د هر
  $w' \in W'$ دپاره: د هر $!A$ دپاره
  $\mSat{M'}{!A}[w']$ هله او يوازې هله چې
  $\mSat{M}{!A}[w']$ وي (دا معنا لري، ځکه چې $W'
  \subseteq W$)۔ په ځانګړي ډول، $\mSat/{M'}{!B}[w]$،
  او~$!B$ په يوه داسې مدل کښې نارښتيا دے چې پر کلي
  چوکاټ ولاړ دے۔
\end{proof}

\begin{figure}[t]
  \centering
  \begin{tikzpicture}[node distance=2cm, auto, thick]
    \clip [rounded corners] (0,0) -- (6,0) -- (6,4) -- (0,4) --  cycle;
    \begin{scope}
      \clip (0,0) -- (2,0)  .. controls (1.5,1.5) and (4.5,1.5)
      .. (4,4) -- (0,4) -- (0,0) -- cycle;
      \filldraw[fill=gray!40] (0,4) -- (0,2) .. controls (1.5,1.5)
      and (4.5,1.5) .. (4.5,0) -- (6,0) -- (6,4) -- (0,4) --  cycle;
    \end{scope}
    \draw [name path=line1] (2,0) .. controls (1.5,1.5) and (4.5,1.5) .. (4,4);
    \draw [name path=line2] (0,2)  .. controls (1.5,1.5) and (4.5,1.5) .. (4.5,0);
    \path [name intersections={of=line1 and line2}];
    \draw [rounded corners, use as bounding box, very thick] (0,0)
    -- (6,0) -- (6,4) -- (0,4) --  cycle; 
    \node at (2,3) {$[w]$} ;
    \node at (1,1) {$[u]$} ;
    \node at (3,0.75) {$[v]$} ;
    \node at (4.75,2) {$[z]$} ;
  \end{tikzpicture}
  \caption{د $W$ په د هم ارزښتۍ ټولګيو وېش۔}
  \ollabel{fig:partition}
\end{figure}

\end{document}
